\chapter{点集拓扑}

18 世纪，Euler 等人发现简单多面体的顶点数 $v$、楞数 $e$、面数 $f$ 总是满足 $v-e+f=2$。
具体地，简单多面体可连续变化成球面。用投影来模拟这个过程：多面体内应存在一种位置，在此处放一盏灯，各顶点和棱都投影在一个外部球面上，这些棱的投影曲线彼此不穿过。比如凸多面体和某些凹多面体，但像厚球壳那样带腔、像甜甜圈那样带洞的多面体是不行的。
这种变化还可形象地视作捏橡皮膜(rubber-sheet)，拓扑学因而又称为\textbf{橡皮几何}。的确，自然界存在直觉上形状相似的几何体，为抽象出共性，可试着“揉”它但不撕裂或粘帖，观察是否能得到另一几何体，如从正方体到球体。这种连续变化间的等价性称作\textbf{同胚}(homeomorphism)。$v-e+f$就是一个多面体在同胚下保持不变的量，即\textbf{拓扑不变量}。
进而我们能将各几何体分门别类，注意力集中于诸如“是否有内外”“是否有洞”“绳上打了几个结”这类问题上的，从而架空距离、面积这些传统几何概念。一个典型例子是，对于理想电路，无论接线长度、方向如何，只要节点不变，电路网络就是等价的。

\begin{figure}[ht]
    \centering
    \includegraphics[width=.8\textwidth]{fig/appx/Mobius.png}
    \caption{M\"obius 环}
\end{figure}

必须指出，用橡皮形变来比喻同胚仍有其局限性。一条纸带不经扭曲，可直接首尾连接得到手环，它具有两个面。扭转半周再相连将得到 M\"obius 环，它只有一个面。扭转一周再相连则又恢复为两个面。这个新环与最初的手环同胚，但却不能通过“揉捏”的方式恢复。如果我们注意到曲面的维度比背景空间的低，而只想关注曲面的内禀性质，则两种手环将视为同⼀空间在 $\R^3$ 内的不同表⽰。所谓的“揉捏”法，实质是 $\R^3$ 的⾃同胚。单看⼆者固然可建⽴同胚，但这种同胚⽆法扩张到整个 $\R^3$ 上。换言之，不存在从 $\R^3$ 的自同胚把两种手环映射在一起，这就解释了问题。这警示我们应更为抽象地看待拓扑。

正如上述例子一样，19 世纪末，人们欲抛掉坐标系而寻求一般集合的拓扑理论。为此要回到同胚的特征：拓扑变化的连续性借助于实数上的 $\epsilon$-$\delta$ 语言，故我们应给出一般集合的连续表述。这一时期可以称为拓扑的分析化，孕育出了\textbf{点集拓扑}。可以说此即最普适的拓扑学。构造这一理论花了很长时间，但我们可简述其关键：极限定义用到了实数集上的小于（序结构）和减法（度量结构），这相当于规定哪类子集是某点邻域、开集和闭集。
因此，拓扑学所研究的对象可以看作脱胎于距离的一种更广义的空间。将这些关键概念抽离出来后，我们就能在拓扑框架下重新得到分析学的定义、命题。现代主流拓扑学正是从这些分析概念出发，一步步走向具象的橡皮几何。
更为深刻的话题将涉及\textbf{代数拓扑}。这一专题将在关于\textbf{上同调}(cohomology)的部分中得以处理，在讨论拓扑分类、积分理论时 de Rham 上同调比较好用。

\section{拓扑空间}

$2^X$ 的子集称为 $X$ 的\textbf{子集族}(collection of subsets)，全体子集族构成 $2^{2^X}$。欲对 $X$ 中每个点挑选⼀族⼦集作为邻域(neighbourhood)，就是给定映射 $\topo N:X\to 2^{2^X}$ 且 $\topo N(x)\ne \mt$。
直观上，邻域只需含有 $x$ 即可，但由于大小任意，可无限接近于 $x$，故有“邻”称。
若 $A\in \topo N(x)$，则 $x$ 为 $A$ 的内点(interior point)，而 $A$ 中全体内点之集 $A^\circ$ 就是其内部(interior)。

希望从分析学中抽取邻域的若干性质，但不能太多，因为要普适到能包括不同于通常定义的情形。新表述不能依赖距离概念。关键在于，注意属于、包含等集合术语，已能描绘距离所表达的含义：$x$ 在其任意邻域内；$x$ 的任意两邻域之交也是 $x$ 的邻域；$x$ 的任意邻域的内部也是其邻域；包含 $x$ 某邻域的集合也是其邻域。
有了邻域，其它概念都好办了：在某全集 $X$ 下，开集(open set)是满足 $A=A^\circ$ 的集合，闭集(closed set)是其补集为开的集合。因此 $\mt,X$ 是既开又闭的。根据 $\epsilon$-$\delta$ 定义，$f:X\to Y$ 在 $x\in X$ 连续，即 $\forall\epsilon>0$ 都 $\exists\delta>0$，使得 $x$ 的 $\delta$ 邻域的像含于 $f(x)$ 的 $\epsilon$ 邻域。

往往直接考察 $f$ 整体的连续性，但邻域毕竟是就某 $x\in X$ 而言的，可预料从开集出发语言将更优雅。非空开集总是某点的某邻域，而邻域的内部必开，故 $f$ 的连续性表述不必出现邻域概念。有两点考虑：$f$ 不一定是满射，但空集可作为 $Y\backslash\Im f$ 的逆像；$f$ 不一定是单射，满足 $f(z)=f(x)$ 的 $z$ 可以很多，但 $f$ 的连续性要求处处连续，因此 $f(x)$ 任意邻域的逆像一定是所有 $z$ 的某邻域之并。综上，$f$ 连续即 $Y$ 中任意开集之逆是 $X$ 的开集。

由于包含邻域之集也是邻域，因此任意开集之并为开；而两邻域之交为邻域，经由归纳法，只能保证有限个开集之交为开；$\mt,X$ 是开的，$x\in X$ 的邻域则是能包含某个含有 $x$ 的开集的集合。这几条性质显然能导出邻域体系，故为等价表述。

\begin{definition}
$\topo{T}\subset 2^X$ 称为 $X$ 的一个\textbf{拓扑}(topology)，若：
\begin{itemize}
    \item $\mt,X \in \topo{T}$；
    \item $\forall \sigma\subset \topo{T}$，$\bigcup_{U\in\sigma} U\in \topo{T}$；
    \item $U,V\in\topo{T}\im U\cap V\in\topo{T}$。或等价地对有限集 $\{A_{i}\} \subset \topo{T}$，$\bigcap_{i} A_{i} \in \topo{T}$。
\end{itemize}
$\topo{T}$ 的元素称为 $X$ 在 $\topo{T}$ 下的\textbf{开子集}，简称\textbf{开集}。
$(X,\topo T)$ 称为 $X$ 关于拓扑 $\topo T$ 的\textbf{拓扑空间}(topological space)。拓扑空间的元素称为\textbf{点}(point)。
\end{definition}

\begin{definition}
    对$x\in X,U\subset X$，若 $\exists O\in\topo{T}$ 使 $x\in O\subset U$，则 $U$ 是 $x$ 的一个\textbf{邻域}。$x$ 的全体邻域构成 $\topo N(x)$。邻域为开则称\textbf{开邻域}，构成 $\topo N(x)\cap\topo T$。
\end{definition}

\begin{theorem}
    $\exists U \in\topo N(x)$ 如何 $\iff \exists$ 开集 $O\ni x$ 如何 $\iff \exists V \in\topo N(x)\cap\topo T$ 如何。
\end{theorem}
\begin{proof}
    $\exists U\in\topo N(x)\im \exists O\in\topo T$ 使 $x\in O\subset U\im \exists$ 开集 $O\ni x\im\exists O\in\topo T$ 使 $x\in O\subset O\im \exists V=O \in\topo N(x)\cap\topo T\im\exists U=V\in\topo N(x)\cap\topo T\subset\topo N(x)$。
\end{proof}
\begin{remark}
    凡命题中出现“点存在邻域…”均可换成“点存在开邻域”“包含该点的开集”。
\end{remark}

%% \setminus or \backslash
\begin{definition}\label{def:closed-set}
    $A\subset X$ 称为\textbf{闭集}，若 $X\backslash A\in\topo{T}$。
\end{definition}

$\topo T$ 的名称若自明或不重要，则拓扑空间常略写为$X$。讨论开集时也常省略全集。De Morgan 律指出：并集之补等于补集之交，因此闭集也能给出拓扑的定义。容易得出等价表述为：$\mt,X$ 是闭的；有限闭集之并为闭；任意闭集之交为闭。

拓扑的选取有很多。若 $\topo T_1\subset\topo T_2$，则称 $\topo T_1$ \textbf{较粗}(coarser)或 $\topo T_2$ \textbf{较细}(finer)。直观上，较细拓扑容纳更多开集，能区分更多点。\textbf{凝聚}(indiscrete)\textbf{拓扑} $\topo T=\{\mt,X\}$ 最粗。\textbf{离散}(discrete)\textbf{拓扑} $\topo T= 2^X$ 最细。
分析学中，$\mt$ 和能表为开区间之并的集合都是 $\R$ 上开集，构成 $\R$ 的\textbf{通常}(usual)\textbf{拓扑} $\topo T_u$。视 $x,y\in\R$ 的距离为 $|x-y|$，则开区间实质是以 $\frac{x+y}{2}$ 为心的一维球。$(\R,\topo T_u)$ 称为\textbf{实直线}。通常拓扑当然比离散、凝聚拓扑更直观，后者一般用于构造反例。

\begin{definition}
    对 $(X,\topo T_X),(Y,\topo T_Y)$，$f: X \to Y$ \textbf{连续}，若 $\forall A\in\topo T_Y$，$f^{-1}[A]\in \topo T_X$。
\end{definition}

\begin{definition}
    双向连续的双射称作\textbf{同胚}。显然同胚时拓扑同势，即开集结构等同。
\end{definition}

\begin{remark}
    同胚的逆也必须连续。
    考虑复平面和 $f(x)=\e^{\i x},x\in [0,2\pi)$。则 $\Im f$ 为单位圆周 $|z|=1$。它是连续双射，但逆不连续。而直观上圆周与区间不能同胚。
\end{remark}

\begin{definition}
    映射 $f:X\to Y$ \textbf{保持}某种性质，就是若 $A\subset X$ 具有这种性质，则 $f[A]$ 也具有。\textbf{开映射}就是保持开的映射。闭映射定义同理。
\end{definition}

\begin{theorem}
    同胚是开映射。
\end{theorem}

\begin{proof}
取任意开集 $U$，同胚 $f:U\to f[U]$。因 $f^{-1}$ 连续，开集 $f^{-1}[f[U]] = U$ 的原像 $f[U]$ 也开。同理同胚也是闭映射。
\end{proof}

\begin{definition}
置 $(X,\topo T)$ 和 $A\subset X$。
\begin{itemize}
    \item \textbf{内部}为开集 $A^\circ:=\bigcup_{O\in\topo T\cap 2^A} O=\{x: \text{$\exists$ 开集 $O\subset A$ 使 $x\in O$}\}$，也即含于 $A$ 的最大开集。
    \item \textbf{闭包}(closure)为闭集 $\bar{A}:=\bigcap_{X\backslash B\in\topo T,A\subset B} B=\{x: \text{$\forall$ 闭集 $B\supset A$ 都 $x\in B$}\}$，也即包含 $A$ 的最小闭集。
    \item $\partial A=\bar{A}\backslash A^\circ$ 称为 $A$ 的\textbf{边界}(boundary)。
\end{itemize}
\end{definition}

\begin{remark}
    显然 $A^\circ\subset A\subset \bar A$。$A^\circ\cup\partial A=\bar A$，$A^\circ\cap\partial A=\mt$。
    
    $\partial A=\mt\iff A^\circ=A=\bar A\iff A$ 既开又闭。
\end{remark}

\begin{theorem}
    $X\backslash A^\circ=\overline{X\backslash A}$，$(X\backslash A)^\circ=X\backslash \bar A$。
\end{theorem}
\begin{proof}
    $A^\circ:=\bigcup_{X\backslash F\in\topo T,X\backslash F\subset A} X\backslash F=X\backslash\bigcap_{X\backslash F\in\topo T,X\backslash A\subset F} F=X\backslash\overline{X\backslash A}$。
    替换即可。
\end{proof}

\begin{theorem}
    $A\in\topo T\eqto A=A^\circ\eqto\forall x\in A,A\in\topo N(x)\eqto\forall x\in A,\exists U\in\topo N(x)$ 使 $U\subset A$。
\end{theorem}

\begin{proof}
    $A\in\topo T,A\subset A\im A\subset A^\circ\iff A=A^\circ$；反之 $A=A^\circ\in\topo T$。

    $A=A^\circ\iff\forall x\in A,\exists O\in\topo T$ 使 $x\in O\subset A\iff\forall x\in A, A\in\topo N(x)\im \forall x\in A,\exists U=A\in\topo N(x)$ 使 $U\subset A\im \forall x\in A,\exists O\in\topo T$ 使 $ x\in O\subset U\subset A \im A=A^\circ$。
\end{proof}

\begin{theorem}
    $A$ 闭 $\iff A=\bar A$
\end{theorem}
\begin{proof}
    $A$ 闭 $\iff X\backslash A=(X\backslash A)^\circ=X\backslash\bar A\iff A=\bar A$。
\end{proof}

\begin{theorem}
    \textbf{内点}：$x\in A^\circ\iff\exists$ 开集 $O\ni x$ 使 $O\subset A^\circ\subset A$；
    \textbf{外点}：$x\notin \bar A\iff\exists$ 开集 $O\ni x$ 使 $O\subset X\backslash\bar A\subset X\backslash A$；
    \textbf{边界点}：$x\in\partial A\iff\forall$ 开集 $O\ni x$ 与 $A,X\backslash A$ 都交。
\end{theorem}
\begin{proof}
    注意 $A^\circ\in\topo T,X\backslash \bar A = (X\backslash A)^\circ\in\topo T$ 即可。并注意 $O\cap A=\mt\iff O\subset X\backslash A$。
\end{proof}

\begin{remark}
    立即可得 $x\in\bar A\iff\forall$ 开集 $O\ni x,O\cap A\ne\mt$。
\end{remark}

\begin{definition}
置 $A\subset X$。\textbf{导集}(derived set)为 $A':=\{x \in X:x\in\overline{A\backslash\{x\}}\}$，其元素称为\textbf{聚点}(accumulation point)或\textbf{极限点}。
$\bar{A}=X$ 时称 $A$ 在 $X$ 中\textbf{稠密}(dense)。
\end{definition}
\begin{remark}
    $x\in A'\eqto \forall U\in\topo N(x),(U\cap A)\backslash\{x\}=U\cap (A\backslash\{x\})\ne\mt$。$U$ 可等价换为任意含 $x$ 开集。
\end{remark}
\begin{example}
    考虑实直线，$0,1/2,1$ 都是 $[0,1)$ 的聚点。
\end{example}

\begin{theorem}
    $\bar A=A\cup A'$。  
\end{theorem}

\begin{proof}
    任给 $x\in A'$。$\forall U\in\topo N(x),(U\cap A)\backslash\{x\}\ne\mt\Rightarrow \forall U\in\topo N(x),U\cap A\ne\mt\eqto x\in\bar A$。说明 $A'\subset \bar A$，此 $A\cup A'\subset\bar A$；
    任给 $x\in\bar A\backslash A'$。$x\in\bar A,x\notin A'\eqto \exists U\in\topo N(x),(U\cap A)\backslash\{x\}=\mt$ 且 $U\cap A\ne\mt\Rightarrow U\cap A=\{x\}\Rightarrow x\in A$。说明 $\bar A\backslash A'\subset A\iff \bar A\subset A\cup A'$。
\end{proof}

\begin{theorem}
    $A$ 闭 $\iff A'\subset A$。
\end{theorem}
\begin{proof}
    $A$ 闭 $\iff A=\bar A\iff A=A\cup A'\iff A'\subset A$。
\end{proof}

\begin{definition}
    设拓扑空间 $X$ 上的映射 $f: X \to Y$，$Y$ 是一个带有零元 $0$ 的集合（$\R,\C$ 或某个拓扑群、线性空间）。那么：
    \begin{itemize}
        \item 零集：$Z(f):=f^{-1}\{0\}=\{x\in X: f(x)=0\}$；
        \item 非零集(cozero set)：$\operatorname{coz}f:=X\backslash Z(f)=\{x\in X: f(x)\ne 0\}$；
        \item \textbf{支集}(support)：$\supp f:=\overline{\operatorname{coz}f}
  = X\backslash (Z(f))^\circ$。
    \end{itemize}
    更方便地，要求 $Y$ 是拓扑空间且 $\{0\}$ 闭（$Y=\R,\C$ 当然满足），从而可谈及连续 $f$。这样 $Z(f)=f^{-1}\{0\}$ 闭，$\operatorname{coz} f$ 开，$\supp f$ 闭。
\end{definition}

\begin{theorem}
     $x\notin \supp f\iff\exists$ 开集 $U\ni x$ 使 $f[U]=\{0\}$。
\end{theorem}
\begin{proof}
    首先 $U\subset Z(f)\iff f[U] \subset \{0\}$。存在元素 $x\in U$，故 $\mt\ne U\subset Z(f)\iff f[U]=\{0\}$，而 $x\in X\backslash\supp f = (Z(f))^\circ \iff \exists$ 开集 $U\ni x$ 使 $U\subset Z(f)$。
\end{proof}
\begin{theorem}
    $\supp f \subset O\subsetneq X\im f[X\backslash O]=\{0\}$。反之需要 $O$ 闭的条件。
\end{theorem}
\begin{remark}
    $\supp f \subset O$ 称为 $f$ 支撑在 $O$ 上。
\end{remark}

\begin{proof}
    $\mt\ne X\backslash O\subset (Z(f))^\circ\subset Z(f)$。
    反之 $\mt \ne X\backslash O\subset Z(f)\xRightarrow{X\backslash O\in\topo T} X\backslash O\subset (Z(f))^\circ$。
\end{proof}

\begin{definition}
    $S:\N\to X,n\mapsto x_n$ 的值域 $\Im S=\left\{x_n\right\}$ 称为\textbf{序列}或\textbf{点列}。
    若 $\exists x \in X$，$\forall U\in\topo N(x),\exists N\in\N$，使 $n>N\im x_n\in U$，则称 $\{x_n\}$ 有\textbf{极限}或\textbf{收敛}于 $x$。
\end{definition}

\begin{remark}
    点列聚点的任意邻域都含点列的无限个点。极限是聚点，但聚点不一定是极限。
\end{remark}

\begin{definition}
    在 $(X,\topo T)$ 下，可定义 $A \subset X$ 的\textbf{相对}(relative)\textbf{拓扑}或\textbf{诱导}(induced)\textbf{拓扑} $\topo T_r=\{O \cap A: O\in\topo T\}$。赋予诱导拓扑的子集称为原集的\textbf{拓扑子空间}(topological subspace)。
\end{definition}

\begin{example}
    对于实直线和 $A=[0,2]$，在 $\topo T_u$ 的诱导下 $B=(1,2]$ 是 $A$ 中的开集。
\end{example}

\begin{definition}
    给定 $(X,\topo T_X),(Y,\topo T_Y)$ 和 $Z=X \times Y$。
    定义
    \[\topo T_Z=\{O\subset Z:A\in\topo T_X,B\in\topo T_Y,\text{$O$ 可表为形如 $A\times B$ 的集合之并}\},\]
    $(Z,\topo T_Z)$ 称为 $X,Y$ 的\textbf{乘积拓扑空间}(product topological space)。
\end{definition}

这与 $\R$ 的通常拓扑类似，都是通过囊括空集和“任意并”操作来生成拓扑。这一方法可总结如下。
\begin{definition}
    置 $(X,\topo T)$。
    $\topo C\subset 2^X$ 称为 $A\subset X$ 的\textbf{覆盖}(cover)，若 $A\subset\bigcup_{O\in\topo C}O$（$A=X$ 时显然取等）。$\topo C\subset\topo T$ 时称\textbf{开覆盖}。
    $\topo S\subset\topo C$ 是 $\topo C$ 的\textbf{子覆盖}(subcover)，若它也是覆盖。
\end{definition}

\begin{definition}
    取 $X$ 的覆盖 $\topo{B}$。若它保持“有限交”性质：对任意 $U_1,U_2\in\topo{B}$，$\exists\topo{F}\subset \topo{B}$ 使 $\bigcup_{U\in\topo{F}}U=U_1\cap U_2$，就称$\topo{B}$ 为\textbf{拓扑基}。可以证明
$\topo{T}=\left\{\bigcup_{U\in\topo{F}}U:\topo{F}\subset\topo{B}\right\}$ 构成拓扑，称为 $\topo{B}$ 的\textbf{生成拓扑}。\textbf{拓扑子基}的元素所有可能的有限交构成拓扑基。
\end{definition}
\begin{theorem}
    $\topo{B}\subset\topo T$，因而是开覆盖。
\end{theorem}

拓扑基是挑出⼀部分开集来代表所有开集，类似地也有一种挑出⼀部分邻域来代表所有邻域的办法。

\begin{definition}
    $x\in X$ 的\textbf{邻域基}或\textbf{局部基} $\topo B_x\subset\topo N(x)$ 满足 $\forall O\in\topo N(x)$，$\exists B\in\topo B_x$ 使 $B\subset O$。
\end{definition}

许多空间往往都难以直接想象或处理，我们可考虑将它分解为简单直观的部分，再重组回去，就能描述复杂的拓扑空间。

\begin{definition}
    对 $(X,\topo T)$ 及等价关系 $\sim$ 规定如下拓扑：设典范投影 $\pi:x\mapsto [x]$，取 $\tilde{\topo T}=\{U \subset \tilde X:\pi^{-1}(U)\in\topo T\}$，称为\textbf{商拓扑}，$(\tilde X,\tilde{\topo T})$ 称为\textbf{商空间}。
\end{definition}
可见，$\pi$ 就好比用来粘合拓扑空间的胶水。
可以证明，设 $\tilde{\topo T}^{\prime}$ 为 $\tilde X$ 上另一个拓扑，且使 $\pi$ 连续，则 $\tilde{\topo T}^{\prime} \subset \tilde{\topo T}$。

\begin{remark}
    \textbf{局部}性质通过各点的邻域描述，但具体要求依性质而定：有时只需每点存在一个具有该性质的邻域，有时则要求这样的邻域构成邻域基。例如，局部可度量只需每点存在一个可度量开邻域，因为度量可限制到更小的子空间；路径连通却未必在缩小邻域后保留，故合理的“局部路径连通”应要求每点存在\emph{任意小}的路径连通开邻域。各术语以下文定义为准。
\end{remark}

\section{可数与分离}
本节目的是给出物理上足够有用的空间，使拓扑性质足够好，不至于太怪异。

在拓扑学中针对一堆研究对象谈它们的可数性时，并不一定是指这些对象的全体可数，因这个要求不是很容易达到。很多时候可数性会定义成：从这堆对象中可以选出可数多个代表来“刻画出”其他所有对象。邻域/拓扑基就是从邻域/开集里面选出来的“代表团”。
和以后要讲的连通性、紧致性不同，直接检查可数性来证明两个空间不同胚的例子并不多见，可数性对于一些数学技巧（特别是数学归纳法）是否能应用至关重要，所以点集拓扑学中有好几个著名定理，如果去掉可数性的前提条件就证明不出来了。用邻域基、拓扑基可定义出两种\textbf{可数定理}。

\begin{definition}\textbf{第一可数}空间的任意点存在可数邻域基，\textbf{第二可数}空间存在可数拓扑基，各记 $C_1,C_2$。字母 C 来⾃于英语“countable”。
\end{definition}

\begin{theorem}
    $C_2$ $\im$ \textbf{Lindel\"of}：任意开覆盖都存在可数子覆盖。
\end{theorem}

\begin{proof}
    取 $X$ 的可数基 $\topo B$ 及任意开覆盖 $\topo U$。置
    \[
        \topo B'=\{B\in\topo B:\exists U\in\topo U,\ B\subset U\}.
    \]
    对每个 $B\in\topo B'$，取 $U_B\in\topo U$ 使 $B\subset U_B$，则 $\topo U'=\{U_B:B\in\topo B'\}$ 至多可数。
    任给 $x\in X$，存在 $U\in\topo U$ 及 $B\in\topo B$ 使 $x\in B\subset U$，故 $B\in\topo B'$，从而 $x\in U_B$。因此 $\topo U'$ 是可数子覆盖。
\end{proof}

\begin{theorem}
    $C_2 \im C_1$。
\end{theorem}

\begin{proof}
    取 $X$ 的可数基 $\topo B$。任给 $x\in X$，置 $\topo B_x=\{B\in\topo B:x\in B\}\subset\topo N(x)$，则 $\topo B_x$ 至多可数。
    对任意 $U\in\topo N(x)$，存在开集 $O$ 使 $x\in O\subset U$，再由基的性质取 $B\in\topo B$ 使 $x\in B\subset O$。故 $B\in\topo B_x$ 且 $B\subset U$，即 $\topo B_x$ 是 $x$ 的可数邻域基。
\end{proof}

再讨论分离。分离意指是否可用拓扑结构（开集、连续映射）区分空间中的两个原本不应重叠的部分。按照对“不重叠两部分”的理解不同、区分方法的不同，有很多种不同判定方法和条件，统称\textbf{分离公理}(separation axiom)。
常用的列举如表 \ref{tab:separation}。字母 T 来⾃于分离公理的德语“das Trennungsaxiom”。
\begin{table}[ht]
    \centering
    \caption{分离公理}
    \label{tab:separation}
    \begin{tabular}{lll}
        \toprule
        符号 & 别称 & 定义\\
        \midrule
        $T_0$ & Kolmogorov & $\forall x \ne y$， $\exists$ 开集 $U$ 使 $(x \in U, y \notin U)\lor (y \in U, x \notin U)$ \\
        $T_1$ & Fréchet & $\forall x \ne y$， $\exists$ 开集 $U, V$ 使 $(x \in U, y \in V)\land (y \notin U , x \notin V)$ \\
        $T_2$ & Hausdorff & $\forall x \ne y$， $\exists$ 不交开集 $U, V$ 使 $x \in U, y \in V$ \\
        $T_3$ & 正则 $^*$ & $T_1$ 并且 $\forall$ 闭集 $F$ 及 $x \notin F$，$\exists$ 不交开集 $U, V$ 使 $x \in U, F \subset V$ \\
        $T_4$ & 正规 $^*$ & $T_1$ 并且 $\forall$ 不交闭集 $A, B$， $\exists$ 不交开集 $U, V$ 使 $A \subset U, B \subset V$ \\
        $T_5$ & 完全正规 $^*$ & 任意子空间都 $T_4$ \\
        \bottomrule
        \multicolumn{3}{l}{\small $^*$ 正则(regular)、正规(normal)有时专指除开 $T_1$ 的第二个条件。}
    \end{tabular}
\end{table}

\begin{theorem}
    $T_1\iff$ 单点集闭。
\end{theorem}
\begin{proof}
    置 $\{x\}$。$\forall y\ne x$ 即 $y\in X\backslash\{x\}$，$y$ 必有开邻域 $V\notni x$ 即 $V\subset X\backslash\{x\}$，故 $X\backslash\{x\}$ 开。
\end{proof}

\begin{theorem}
    $T_5\im T_4\im T_3\xRightarrow{(T_3\Rightarrow T_1)}T_2 \im T_1\im T_0$。
\end{theorem}

\begin{proof}
    $X$ 是自身的子空间，故 $T_5\im T_4$。
    在 $T_4$ 空间中，单点集闭，故对任意闭集 $F$ 及 $x\notin F$，可用不交开集分别包含 $\{x\},F$，即 $T_3$。
    在 $T_3$ 空间中，任给 $x\ne y$，因 $\{y\}$ 闭，可用不交开集分别包含 $x,y$，即 $T_2$。
    不交的开邻域各不含另一点，故 $T_2\im T_1$；$T_1$ 中取其中一个开邻域即得 $T_0$。
\end{proof}

\begin{example}
    考虑 $\R^2$ 上直线 $y=0,1$。取两条直线之并，取等价关系 $(x,0) \sim(x,1),x>0$ 来构建商拓扑空间，也就是把两条线的 $x>0$ 部分黏合一起。$(0,0),(0,1)$ 的任意开邻域存在黏合部分 $x>0$，也就没有不交开邻域，因而不是 $T_2$ 空间。
\end{example}

\begin{example}
    给 $\N_+$ 定义拓扑，开集取 $\mt,\N_+$ 以及 $\{1,2,3,\cdots,n\}$。这个空间既非 $T_2$ 亦非紧致的。
\end{example}

\begin{theorem}[收缩(shrink)]
    设 $X$ 为 $T_1$ 空间，则显然：

    $T_3\iff$ 对任意点 $x\in X$ 及其任意邻域 $O$，存在开集 $U$ 使
$x\in U, \bar U \subset O$。

$T_4\iff$ 对任意闭集 $F\subset X$ 与任意开集 $O\supset F$，存在开集 $U$ 使 $F\subset U, \bar U \subset O$。
\end{theorem}

以 Hausdorff 拓扑空间为对象、连续映射为态射可构建\textbf{拓扑空间范畴} $\cate{Top}$。合成与恒等态射的定义同理。
类似地定义带基点的拓扑空间范畴 $\cate{Top}_\bullet$. 我们也希望赋予同态集 $\hom(X,Y)$ 额外的结构，例如紧开拓扑，使得 $\hom(Y, Z) \times \hom(X, Y) \to \hom(X, Z)$ 成为连续映射；我们更希望能有自然的同构
\[ \hom(X \times Y, Z) \rightiso \hom(X, \hom(Y, Z)). \]
相关的点集拓扑问题颇为棘手，为了在确保良好的范畴性质的同时容许充分广的拓扑空间，在同伦论里一般选用 $\cate{Top}$ 的一个子范畴 $\cate{CGHaus}$，称为\textbf{紧生成 Hausdorff 空间范畴}.

\section{度量空间}

我们看到，$\R$ 的距离便可生成一个拓扑基。距离概念的推广自然是在任意集合中，给定满足三角不等式、正定性的对称实函数。开区间也可相应地推广。实际上，具有应用意义的拓扑空间大多具有距离概念。

\begin{definition}\label{def:metric_topo}
    置集合 $X$ 和\textbf{距离}（或\textbf{度量}）$d: X \times X \to [0,\infty)$，$\forall x,y,z\in X$ 满足：
    \begin{itemize}
        \item $d(x,y)=d(y,x)$；
        \item $d(x,y) = 0\eqto x=y$；
        \item $d(x,z) \leqslant d(x,y)+d(y,z)$。（三角不等式）
    \end{itemize}
    $B(x,r)=\{y \in X: d(x, y)<r\}$ 称为以 $x \in X$ 为中心、$r$ 为半径的\textbf{开球}(open ball)。
    全体开球 $\{B(x,\epsilon):x\in X,\epsilon >0\}$ 就是一个拓扑基，其生成拓扑 $\topo T_d$ 就是 $(X,d)$ 的\textbf{度量诱导拓扑}。$(X,d,\topo T_d)$ 称为\textbf{度量空间}(metric space)。存在度量 $d$ 生成的拓扑，称\textbf{可度量}。
\end{definition}

\begin{remark}
    此处的距离须非负，与双线性函数意义下的度规不同。不定度规的非退化性并不排除非零矢量的模方为零，不能直接将其模长用作距离。
\end{remark}

\begin{definition}
    \textbf{局部可度量}就是 $\forall x\in X$ 都存在开集 $O\ni x$ 可度量。
\end{definition}

\begin{example}
    \textbf{离散度量}满足 $d_{0}(x,y)=1$ 若 $x \ne y$。这是一个描述离散拓扑的办法。
\end{example}

\begin{definition}
    置度量空间 $X$。度量 $d_{1},d_{2}$ 是\textbf{等价度量}，若 $\exists a,b>0$，使对 $\forall x,y \in X$，有 $a d_{1}(x, y) \leqslant d_{2}(x, y) \leqslant b d_{1}(x, y)$。
\end{definition}

\begin{remark}
    显然等价度量诱导出相同拓扑。
\end{remark}

\begin{example}
    置 $x=\left(x^1,\cdots,x^{n}\right),y=\left(y^{1},\cdots,y^{n}\right)\in\R^{n}$。$p$-度量为
\[
d_{p}(x, y):=\left(\sum_{k=1}^{n}\left|x^{k}-y^{k}\right|^{p}\right)^{1 / p},\quad p \geqslant 1.
\]
常取$p=2$，称为\textbf{欧氏度量}。
$\R$ 的开球即开区间。$\R^2$ 的开球称为\textbf{开圆盘}。极限为
\[
d_{\infty}(x,y):=\lim_{p\to\infty}d_{p}(x, y)=\max _{1 \leqslant k \leqslant n}\left\{\left|x^{k}-y^{k}\right|\right\}.
\]
任意 $p$-度量等价，诱导出通常拓扑 $\topo T_u$。当然，亦可视作 $\R$ 通常拓扑的乘积拓扑。
\end{example}

\begin{theorem}
    可度量 $\im T_5$。
\end{theorem}

\begin{proof}
    子空间 $Y\subset X$ 上的限制度量满足 $B_Y(y,r)=Y\cap B_X(y,r)$，故诱导相对拓扑。因此只需证度量空间是 $T_4$ 的。
    单点集显然闭。任给不交闭集 $A,B$，对每个 $a\in A,b\in B$，取 $r_a,s_b>0$ 使 $B(a,2r_a)\cap B=\mt$，$B(b,2s_b)\cap A=\mt$。置开集
    \[
        U=\bigcup_{a\in A}B(a,r_a),\qquad V=\bigcup_{b\in B}B(b,s_b),
    \]
    则 $A\subset U,B\subset V$。若 $z\in U\cap V$，则存在 $a\in A,b\in B$ 使
    \[
        d(a,b)\leqslant d(a,z)+d(z,b)<r_a+s_b\leqslant 2\max\{r_a,s_b\}\leqslant d(a,b),
    \]
    矛盾。故 $U\cap V=\mt$，即得 $T_4$。
\end{proof}

我们已悉知 $\epsilon$-$\delta$ 语言可视为连续性在通常拓扑的特例。现在再用度量空间重新叙述。
置 $(X,d_X,\topo T_X),(Y,d_Y,\topo T_Y)$ 和映射 $f:X\to Y$。$f$ 在 $x \in X$ 处对度量诱导拓扑连续，当且仅当 $\forall\epsilon>0,\exists\delta>0$ 使 $B(x,\delta) \subset f^{-1} [B(f(x),\epsilon)]$。
即 $f^{-1} [B(f(x),\epsilon)]$ 是 $x$ 点邻域。
换言之，$d_X(x,x^{\prime})<\delta\Rightarrow d_Y(f(x),f(x^{\prime}))<\epsilon$。
若 $f$ 在任意点连续，即 $\forall x \in X,\epsilon > 0$，$f^{-1} [B(f(x),\epsilon)]\in\topo T_X$，结合“任意并”性质和逆像的保并性，可知 $\forall U\in\topo T_Y$，$f^{-1}[U]\in\topo T_X$，即 $f$ 连续。

\begin{theorem}
    $\R^n$ 开球和 $\R^n$ 同胚。
\end{theorem}
\begin{proof}
    不失一般性，考虑单位开球 $B(0,1)$。置 $f:B(0,1)\to \R^n$ 为
\eq{\label{eq:B-Rn}
f(x)=\frac{x}{1-|x|},
\quad
f^{-1}(y)=\frac{y}{1+|y|}.
}
它是双向连续的双射。
\end{proof}

对于度量空间，点列 $\left\{x_n\right\}$ 收敛于 $x$ 等价于 $\lim_{n\to\infty}d(x, x_n)=0$。
点列称为基本列或 Cauchy 列，即$\forall\epsilon>0,\exists N\in\N$使 $k,l>N\Rightarrow d(x_{k},x_{l})<\epsilon$。
度量空间称为\textbf{完备的}(complete)，若其中任意 Cauchy 列都收敛于其中的点。
%
实直线上闭区间完备，但开区间不完备：只需注意 $\{1/n\}_{n=2}^\infty$ 在 $(0,1)$ 上无极限。
$\R^{n}$ 当然完备。
一致连续性即 $\forall \epsilon>0,\exists\delta>0$ 
使 $\forall x_{1},x_{2} \in X$，$d(x_{1},x_{2})<\delta\Rightarrow d(f(x_{1}),f(x_{2}))<\epsilon$。
注意一致连续要更强：对整个空间要求一致的 $\delta$。

\begin{definition}
    度量空间 $X$ 的子集称为\textbf{有界的}，若存在开球包含它。否则称\textbf{无界的}。
\end{definition}

\begin{remark}
    完备性、有界性都不是拓扑不变性。例如，$(0,1)$ 与实直线同胚，但前者不完备、有界，后者完备、无界。
\end{remark}

置度量空间 $X$。$f: X \to X$ 称为（严格）\textbf{压缩映射}(contraction)，若 $\exists c\in(0,1)$，使 $\forall x,y \in X$ 且 $x\ne y$，有 $d(f(x),f(y)) \leqslant c d(x,y)$。
仿照分析学可类似证明 \textbf{Banach 不动点定理}：若度量空间 $X$ 完备，$f$ 为一个压缩映射，则 $f(x)=x$ 的解存在且唯一。这是非线性泛函分析的一个有用工具。

%%%

\begin{definition}
    \textbf{有限 $\epsilon$-网}(net)是有限集 $\{x_n\}\subset X$，使 $\forall x\in X$，$\exists x_i$ 使 $d(x,x_i)<\epsilon$。换言之，$X\subset\bigcup_{i=1}^n B(x_i,\epsilon)$。
\end{definition}

\begin{definition}
    $(X,d)$ 称为\textbf{全有界的}(totally bounded)，若对任意 $\epsilon>0$，$X$ 都存在有限 $\epsilon$-网。
\end{definition}

\begin{remark}
    全有界 $\im$ 有界。取 $\epsilon=1$，由三角不等式即得。若有限 $1$-网为 $\{x_1,\ldots,x_n\}$，则 $X\subset B (x_1,\,1+\max_{1\le i\le n}d(x_1,x_i) )$，故有界。
\end{remark}

\begin{example}
    无穷维 Hilbert 空间 $\ell^2$ 完备，但不全有界。标准基矢满足 $\|e_n-e_m\|=\sqrt2$（$n\ne m$），故半径小于 $\sqrt2/2$ 的球至多含一个 $e_n$，有限个这样的球不能覆盖 $\ell^2$。
\end{example}

%%%

\section{紧致性}

分析学给出了若干等价的实数定理，其中一条是著名的\textbf{Heine-Borel 定理}：若有界闭集能被某开集族覆盖，则可从中找到有限个开集，构成依旧能覆盖的子族。
19 世纪中叶，Dirichlet 在试图严格化函数连续性时，率先用有限覆盖技术，证明了闭区间上的连续实函数一致连续。
不过他的结果半个世纪后才发表，期间 Heine, Weierstrass 等人各自独立地用类似技术也证明了该结论。Borel 利用这些技术证明了上述的有限覆盖定理。后来 Lebesgue 等人才把它推广到了一般情形。
有界性是与距离密切相关的，但有限覆盖性只涉及开集。可见关于有界闭集的分析学命题，很多都能延伸到点集拓扑。

\begin{definition}
    $A$ 称\textbf{紧致的}(compact)或\textbf{紧的}，若其任意开覆盖都存在有限子覆盖。
\end{definition}

\begin{example}
    实直线上，单点集 $\{x\}\subset X$ 必紧。$\R$、开区间、半开区间非紧。
\end{example}
\begin{proof}
    对 $\{x\}$ 任取开覆盖 $\topo C$，则 $\exists O\in\topo C$ 使 $x\in O$，$\{O\}\subset\topo C$ 是有限集。
    不妨以 $(0,1]$ 为例，存在开覆盖 $\{(1/n,2):n\in\N\}$ 不具有限子覆盖；其余情形同理。
\end{proof}

\begin{definition}
    若 $K\subset X$ 中任何序列都包含收敛子列，则称 $K$ 是\textbf{相对紧的}。若这些子列还收敛于 $K$ 中，则称\textbf{自列紧的}，简称\textbf{列紧}\footnote{列紧在部分书中指相对紧。}。
\end{definition}

%%%
\begin{remark}
    一般拓扑空间中，紧与列紧互不蕴含。
\end{remark}

\begin{theorem}\label{thm:compactness}
    对度量空间，紧 $\iff$ 列紧 $\iff$ 全有界 $\land$ 完备。
\end{theorem}
\begin{proof}
    紧 $\im$ 列紧：任给点列 $\{x_n\}$。若每点都有仅含有限项的开邻域，则有限子覆盖将使点列仅有有限项，矛盾。故存在 $x\in X$，其任意邻域都含无限项；逐次取 $n_k>n_{k-1}$ 使 $x_{n_k}\in B(x,1/k)$，即得收敛子列。

    列紧 $\im$ 全有界且完备：若不全有界，则存在 $\epsilon>0$，可逐次取 $x_n\notin\bigcup_{i<n}B(x_i,\epsilon)$。各项间距离不小于 $\epsilon$，故无收敛子列，矛盾。又任意 Cauchy 列都有收敛子列，因而本身收敛，故完备。

    全有界且完备 $\im$ 紧：反设开覆盖 $\topo C$ 无有限子覆盖，置 $A_0=X$。对每个 $n\geqslant1$，用有限个半径 $2^{-n}$ 的开球覆盖 $X$，其中必有一个与 $A_{n-1}$ 的交 $A_n$ 仍不能被 $\topo C$ 中有限个开集覆盖。取 $x_n\in A_n$；由 $A_{n+1}\subset A_n$，当 $m\geqslant n$ 时 $d(x_m,x_n)<2^{1-n}$，故 $\{x_n\}$ 为 Cauchy 列，收敛于某 $x\in X$。
    取 $U\in\topo C$ 使 $x\in U$。任给 $y\in A_n$，当 $m\geqslant n$ 时 $d(y,x_m)<2^{1-n}$；令 $m\to\infty$ 得 $d(y,x)\leqslant2^{1-n}$，故充分大 $n$ 时 $A_n\subset U$，矛盾。
\end{proof}

\begin{remark}
    完备子空间在任何度量空间中都闭，故无需另加闭性。若环境空间 $X$ 完备，则 $K\subset X$ 闭 $\iff K$ 完备，从而 $K$ 紧 $\iff K$ 全有界且在 $X$ 中闭。
\end{remark}

\begin{example}[Bolzano-Weierstrass]
    实直线上，列紧 $\iff$ 有界 $\land$ 闭。
\end{example}
\begin{proof}
    $K\subset\R$ 有界时，将包含它的有界区间分成有限段，每段长度小于 $\epsilon$，从每段与 $K$ 的非空交中取一点，即得有限 $\epsilon$-网。结合全有界 $\im$ 有界，知 $K$ 有界 $\iff$ 全有界。
    $\R$ 完备，故闭子集完备；反之，若 $K$ 完备，则任意 $x\in\bar K$ 都是 $K$ 中某点列的极限，该列为 Cauchy 列，故 $x\in K$。于是 $K$ 闭 $\iff$ 完备，由定理 \ref{thm:compactness} 即得。
\end{proof}

\begin{example}[Heine-Borel]
    结合 \ref{thm:compactness} 知，实直线上，紧 $\iff$ 有界 $\land$ 闭。此外也可直接由紧性证明有界性：若 $K\subset\R$ 紧，则开覆盖 $K\subset\bigcup_{n\in\N^+}(-n,n)=\R$ 有有限子覆盖，即存在 $N = \max_i{n_i}\in\N^+$ 使 $K\subset \bigcup_{i=1}^m(-n_i,n_i)=(-N,N)$，即有界。
\end{example}
%%%

\begin{definition}
    \textbf{局部紧}是任意点都存在\textbf{邻域}相对紧（闭包紧）。
\end{definition}

紧致要求只用有限的开集就能覆盖全集，但这太过严苛。
可以尝试弱化，即全局上仍可能需无穷个开集，但要求每个点附近只相交有限个开集。
这样无穷和在每点附近仍可转化为对有限项的拼接，便于操作。

%%%
\begin{definition}
置 $X$ 的覆盖 $\topo V,\topo U$，若 $\forall V\in\topo V$ 都 $\exists U\in\topo U$ 使 $V\subset U$，
则称 $\topo V$ 是 $\topo U$ 的一个\textbf{细化}或\textbf{加细}(refinement)。默认讨论开覆盖，故称\textbf{开加细}。
\end{definition}

\begin{definition}
覆盖 $\topo V$ 称\textbf{局部有限}(locally finite)，若 $\forall x\in X$，$\exists$ $x$ 的邻域只与 $\topo V$ 中有限多个集合相交：$|\{V\in\topo V:\ O\cap V\ne\varnothing, O\in\topo N(x)\}|<\aleph_0$。
\end{definition}


\begin{definition}
    \textbf{仿紧}或\textbf{副紧}(paracompact)的拓扑空间中，任意开覆盖都存在局部有限的开加细。
\end{definition}
%%%
\begin{remark}
    更原始的定义是称“任意开覆盖都存在开星加细”：称 $\topo V$ 是覆盖 $\topo U$ 的星加细，对每个 $V\in\topo V$，存在 $U\in\topo U$ 使 $\operatorname{St}(V,\topo V):=\bigcup\{O\in\topo V:\ O\cap V\ne\varnothing\}\subset U$。
    但“局部有限的开加细”最实用。
\end{remark}

\begin{theorem}[Stone]
    可度量 $\im$ 仿紧。
\end{theorem}
\begin{proof}
    % M. E. Rudin, A new proof that metric spaces are paracompact,
    % Proc. Amer. Math. Soc. 20 (1969), 603.
    将任意开覆盖 $\{U_\alpha\}$ 的指标集良序化，置 $a(x)=\min\{\alpha:x\in U_\alpha\}$，$r_n=2^{-n}$。依 $n\geqslant1$ 递归定义 $V_{\alpha n}$ 为所有满足
    \[
        a(x)=\alpha,\qquad
        x\notin\bigcup_{m<n,\,\beta}V_{\beta m},\qquad
        B(x,3r_n)\subset U_\alpha
    \]
    的开球 $B(x,r_n)$ 之并。显然 $V_{\alpha n}\subset U_\alpha$。任给 $x\in X$，取 $n$ 使 $B(x,3r_n)\subset U_{a(x)}$；若 $x$ 未被前 $n-1$ 层覆盖，则 $x\in V_{a(x)n}$。故 $\{V_{\alpha n}\}$ 是开加细。

    固定 $n$。若 $\alpha<\beta$，分别取构成 $V_{\alpha n},V_{\beta n}$ 的球心 $y,z$。由 $a(z)=\beta$ 知 $z\notin U_\alpha$，而 $B(y,3r_n)\subset U_\alpha$，故 $d(y,z)\geqslant3r_n$。由三角不等式，两球内任意两点的距离大于 $r_n$，故不同 $V_{\alpha n}$ 中的点也相距大于 $r_n$。

    任给 $p\in X$，取 $p\in V_{\alpha k}$，再取 $N>k$ 使 $B(p,2r_N)\subset V_{\alpha k}$。当 $n\geqslant N$ 时，第 $n$ 层的球心均在 $V_{\alpha k}$ 外，且 $r_n\leqslant r_N$，故该层不与 $B(p,r_N)$ 相交。当 $n<N$ 时，$r_n\geqslant2r_N$，由上述距离估计，$B(p,r_N)$ 至多与该层的一个集合相交。因此 $\{V_{\alpha n}\}$ 局部有限。
\end{proof}

\section{度量化定理*}

\begin{theorem}
    $T_2\xRightarrow{\text{仿紧}} T_4$ 
\end{theorem}
\begin{proof}
    局部有限族满足 $\overline{\bigcup_i V_i}=\bigcup_i\bar V_i$，因每点附近只涉及有限并。

    先证 $T_3$。任给闭集 $F$ 及 $x\notin F$，由 $T_2$，对每个 $y\in F$ 可取开邻域 $U_y$ 使 $x\notin\bar U_y$。取开覆盖 $\{X\backslash F\}\cup\{U_y:y\in F\}$ 的局部有限开加细 $\topo V$，置
    \[
        W=\bigcup_{V\in\topo V,\,V\cap F\ne\mt}V.
    \]
    每个参与并集的 $V$ 必含于某个 $U_y$，故 $F\subset W$ 且 $x\notin\bar W$。于是 $W,X\backslash\bar W$ 分离 $F,x$，即得 $T_3$。

    再取不交闭集 $A,B$。由 $T_3$，对每个 $a\in A$ 可取开邻域 $O_a$ 使 $\bar O_a\cap B=\mt$。取 $\{X\backslash A\}\cup\{O_a:a\in A\}$ 的局部有限开加细，令 $U$ 为其中与 $A$ 相交者之并。同理 $A\subset U$ 且 $\bar U\cap B=\mt$，故 $U,X\backslash\bar U$ 分离 $A,B$，即得 $T_4$。
\end{proof}

\begin{lemma}[Urysohn]\label{lem:urysohn-separation}
    $T_4$ 空间中，对任意不交闭集 $A,B$，存在连续函数 $f:X\to[0,1]$ 使 $f|_A=0,f|_B=1$（从而将二者分离）。
\end{lemma}
\begin{proof}
    置 $D=\{k/2^n:n,k\in\N,\ 0\leqslant k\leqslant2^n\}$，$O_1=X\backslash B$。由收缩定理取 $A\subset O_0\subset\bar O_0\subset O_1$，再逐次插入二等分点对应的开集，使 $r<s\im\bar O_r\subset O_s$。令
    \[
        f(x)=\inf\bigl(\{r\in D:x\in O_r\}\cup\{1\}\bigr).
    \]
    则 $f|_A=0,f|_B=1$。对 $0\leqslant t\leqslant1$，
    \[
        \{f<t\}=\bigcup_{r<t}O_r,\qquad
        \{f>t\}=\bigcup_{r>t}(X\backslash\bar O_r),
    \]
    右端均开，故 $f$ 连续。
\end{proof}

Urysohn 引理用连续函数分离闭集。进一步，闭子集上已给定的连续实值函数，也可延拓到全空间。

\begin{theorem}[Tietze 延拓定理]\label{thm:tietze-extension}
    设 $X$ 为 $T_4$ 空间，$A\subset X$ 闭，则任意连续函数 $f:A\to\R$ 都有连续延拓 $F:X\to\R$，即 $F|_A=f$。若 $f$ 取值于闭区间 $[a,b]$，则可使 $F$ 仍取值于 $[a,b]$。
\end{theorem}
\begin{proof}
    先证有界情形，经线性变换只须考虑 $f:A\to[-1,1]$。任给连续函数 $g:A\to[-r,r]$（$r>0$），集合
    \[
        A_-:=\{x\in A:g(x)\leqslant-r/3\},\qquad
        A_+:=\{x\in A:g(x)\geqslant r/3\}
    \]
    是 $X$ 中的不交闭集。由 Urysohn 引理，存在连续函数 $h:X\to[-r/3,r/3]$，在 $A_-,A_+$ 上分别为 $-r/3,r/3$，故 $|g-h|\leqslant2r/3$ 在 $A$ 上成立。

    从 $g=f,r=1$ 开始，对每次的余项重复此构造，得到连续函数 $h_n:X\to\R$，满足
    \[
        |h_n|\leqslant\frac13\left(\frac23\right)^{n-1},\qquad
        \left|f-\sum_{k=1}^nh_k\right|\leqslant\left(\frac23\right)^n\quad\text{在 }A\text{ 上}.
    \]
    因此 $F:=\sum_{n=1}^\infty h_n$ 一致收敛，连续且取值于 $[-1,1]$，并有 $F|_A=f$。一般闭区间由线性变换得到；$a=b$ 时取常函数即可。

    对无界的 $f$，令 $u=f/(1+|f|):A\to(-1,1)$。由已证情形取连续延拓 $H:X\to[-1,1]$。闭集 $B:=\{x:|H(x)|=1\}$ 与 $A$ 不交，故由 Urysohn 引理可取连续函数 $\chi:X\to[0,1]$，在 $A$ 上为 $1$、在 $B$ 上为 $0$。于是 $G:=\chi H$ 处处满足 $|G|<1$，且 $G|_A=u$，所求延拓即
    \[
        F=\frac{G}{1-|G|}.\qedhere
    \]
\end{proof}
\begin{remark}
    闭性不可略去。例如 $f(x)=1/x$ 在 $(0,1)\subset\R$ 上连续，却不能连续延拓到 $\R$。定理不要求 $f$ 有界，也不保证延拓唯一。
\end{remark}

\begin{theorem}[Urysohn-Lindel\"of]\label{thm:urysohn-metrization}
    $T_3\xLeftrightarrow{C_2} T_4\xLeftrightarrow{C_2}$ 可度量
\end{theorem}
\begin{proof}
    先证 $C_2\land T_3\im T_4$。任给不交闭集 $A,B$，由可数基及收缩定理，取覆盖 $A,B$ 的可数开集族 $\{U_n\},\{V_n\}$，使 $\bar U_n\cap B=\bar V_n\cap A=\mt$。置
    \[
        U=\bigcup_{n\geqslant1}\left(U_n\backslash\bigcup_{j\leqslant n}\bar V_j\right),\qquad
        V=\bigcup_{n\geqslant1}\left(V_n\backslash\bigcup_{j\leqslant n}\bar U_j\right).
    \]
    则 $U,V$ 为分别包含 $A,B$ 的不交开集，故 $T_4$。

    再设 $C_2\land T_4$，取可数基 $\{B_i\}$。对每对 $\bar B_i\subset B_j$，由引理 \ref{lem:urysohn-separation} 取连续函数 $f_{ij}:X\to[0,1]$，在 $\bar B_i$ 上为 $1$，在 $X\backslash B_j$ 上为 $0$。将这些函数列为 $\{f_n\}_{n\geqslant1}$（有限时补零），置
    \[
        d(x,y)=\sum_{n=1}^\infty2^{-n}|f_n(x)-f_n(y)|.
    \]
    任给开集 $O\ni x$，由收缩定理可取 $x\in B_i\subset\bar B_i\subset B_j\subset O$，故某个 $f_n$ 在 $x$ 处为 $1$、在 $X\backslash O$ 上为 $0$。此族因而分离点，$d$ 是度量，且 $B_d(x,2^{-n})\subset O$。另一方面，级数一致收敛，故 $d(x,\cdot)$ 连续，度量开球在原拓扑中开。因此 $d$ 诱导原拓扑。逆向蕴涵已由前文证明。
\end{proof}

\begin{theorem}[Smirnov]
    $\text{仿紧}\land T_2\land\text{局部可度量}\im$ 可度量
\end{theorem}
\begin{proof}
    由前文知 $X$ 为 $T_4$。取由非空可度量开集组成的局部有限覆盖 $\{U_i\}$，各带有相容度量 $d_i\leqslant1$（将原度量换为 $\min\{d_i,1\}$ 即可）。先用收缩定理，再取局部有限开加细，得到开覆盖 $\{V_j\}$ 及 $\bar V_j\subset U_{i(j)}$。置 $F_i=\bigcup_{i(j)=i}\bar V_j$；由局部有限性，$F_i$ 闭，且 $F_i\subset U_i$、$\bigcup_iF_i=X$。
    由引理 \ref{lem:urysohn-separation}，取连续函数 $\phi_i:X\to[0,1]$，在 $F_i$ 上为 $1$，在 $X\backslash U_i$ 上为 $0$。对 $z\in U_i$ 置
    \[
        g_{i,z}(x)=
        \begin{cases}
            \phi_i(x)d_i(x,z),&x\in U_i,\\
            0,&x\notin U_i,
        \end{cases}
    \]
    并令
    \[
        \begin{aligned}
            p_i(x,y)&=|\phi_i(x)-\phi_i(y)|+\sup_{z\in U_i}|g_{i,z}(x)-g_{i,z}(y)|,\\
            d(x,y)&=\sum_i p_i(x,y).
        \end{aligned}
    \]
    各 $p_i$ 非负、对称且满足三角不等式，和式由局部有限性只有有限项。若 $d(x,y)=0$，取 $x\in F_i$，则 $\phi_i(y)=1$，再取 $z=x$ 得 $d_i(x,y)=0$，故 $x=y$。因此 $d$ 是度量。

    当 $y\to x$ 时，若 $x\in U_i$，则在 $U_i$ 内有 $p_i(x,y)\leqslant2|\phi_i(x)-\phi_i(y)|+d_i(x,y)\to0$；若 $x\notin U_i$，则 $p_i(x,y)\leqslant2\phi_i(y)\to0$。结合局部有限性得 $d(x,y)\to0$。
    反之，任给开集 $O\ni x$，取 $x\in F_i$ 及 $0<\epsilon<1/2$ 使 $B_{d_i}(x,2\epsilon)\subset O\cap U_i$。若 $d(x,y)<\epsilon$，则 $\phi_i(y)>1/2$，且取 $z=x$ 得 $\phi_i(y)d_i(x,y)\leqslant d(x,y)<\epsilon$，故 $y\in O$。因此 $d$ 诱导原拓扑。
\end{proof}

\begin{theorem}[Nagata-Smirnov]
    可度量 $\iff T_3\land$ 存在 $\sigma$-局部有限基，即拓扑基可写为 $\topo B=\bigcup_{n=1}^\infty\topo B_n$，其中每个 $\topo B_n$ 都局部有限。
\end{theorem}
\begin{remark}
    这是另⼀个经典研究成果。证明毕竟精巧，实际应⽤中也没有很多可以模仿着去做的机会，读者粗读即可。
    这里可数的是层数，各层本身不必可数。
\end{remark}
\begin{proof}
    若 $X$ 可度量，则已知 $T_3$。由 Stone 定理，对每个 $n\geqslant1$，取开球覆盖 $\{B_d(x,2^{-n}):x\in X\}$ 的局部有限开加细 $\topo B_n$。任给开集 $O\ni x$，取 $n$ 使 $B_d(x,2^{1-n})\subset O$；含 $x$ 的某个 $B\in\topo B_n$ 含于一个半径 $2^{-n}$ 的球，故 $B\subset O$。因此 $\bigcup_n\topo B_n$ 是所求基。

    反之，设 $X$ 为 $T_3$ 且有上述基。对开集 $O$，置
    \[
        W_n(O)=\bigcup\{B\in\topo B_n:\bar B\subset O\}.
    \]
    局部有限性使 $\overline{W_n(O)}\subset O$，而收缩定理及基的性质给出 $O=\bigcup_nW_n(O)$。对不交闭集 $A,B$，分别取 $W_n(X\backslash B),W_n(X\backslash A)$，沿用定理 \ref{thm:urysohn-metrization} 证明中的可数开集族分离构造，即得 $T_4$。

    现在用连续函数表示每个基元素 $B$：在 $B$ 内为正，在外为零。由 Urysohn 引理，对每个 $m$ 取连续函数 $h_{Bm}:X\to[0,1]$，在 $\overline{W_m(B)}$ 上为 $1$，在 $X\backslash B$ 上为 $0$。则
    \[
        f_B=\sum_{m=1}^\infty2^{-m}h_{Bm}
    \]
    一致收敛，故 $f_B$ 连续，且 $0\leqslant f_B\leqslant1$、$\{f_B>0\}=B$。

    同层取函数值之差的上确界，再将各层加权相加：
    \[
        p_n(x,y)=\sup_{B\in\topo B_n}|f_B(x)-f_B(y)|,\qquad
        d(x,y)=\sum_{n=1}^\infty2^{-n}p_n(x,y),
    \]
    空层约定 $p_n=0$。各 $p_n$ 非负、对称且满足三角不等式；基能分离点，故 $d(x,y)=0\im x=y$。因此 $d$ 是度量。

    局部有限性使 $p_n(x,\cdot)$ 在每点附近只涉及有限个 $f_B$，故连续；又 $p_n\leqslant1$，级数一致收敛，故 $d(x,\cdot)$ 连续。反之，任给开集 $O\ni x$，取 $x\in B\subset O$，其中 $B\in\topo B_n$。因 $f_B(x)>0$ 且 $f_B$ 在 $B$ 外为零，有
    \[
        B_d\bigl(x,2^{-n}f_B(x)\bigr)\subset B\subset O.
    \]
    因此 $d$ 诱导原拓扑。
\end{proof}

\begin{theorem}
    $T_2 \xRightarrow{\text{局部紧}} T_3$
\end{theorem}
\begin{proof}
    任给闭集 $F$ 及 $x\notin F$。由局部紧性，取 $x$ 的紧闭邻域 $K$ 及开集 $O$ 使 $x\in O\subset K$。
    对每个 $y\in F\cap K$，由 $T_2$ 取不交开集 $U_y\ni x,V_y\ni y$。因 $F\cap K$ 紧，可取有限个 $V_{y_1},\ldots,V_{y_n}$ 覆盖它。置
    \[
        U=O\cap\bigcap_{j=1}^nU_{y_j},\qquad
        V=(X\backslash K)\cup\bigcup_{j=1}^nV_{y_j}.
    \]
    则 $U,V$ 为不交开集，且 $x\in U,F\subset V$。又 $T_2\im T_1$，故 $T_3$。
\end{proof}

\begin{definition}
    置开覆盖 $\topo P=\{O_i\}_{i\in I}$。若函数族 $\{\rho_i\}_{i\in I}$ 满足
    \begin{itemize}
        \item $0\leqslant\rho_i\leqslant 1$；
        \item $\rho_i$ 支撑在 $O_i$ 上；
        \item $\sum_i \rho_i=1$,
    \end{itemize}
    则称从属于 $\topo P$ 的一个\textbf{单位分解}(partition of unity)。若 $\rho_i$ 连续则称\textbf{连续单位分解}。
\end{definition}

\begin{remark}
    指标集 $I$ 可以是任意集合，有限、可数或不可数均可，定义本身不要求可数。在下述仿紧 $T_2$ 条件下，若另有 Lindel\"of 性（例如第二可数），则可先从开覆盖抽取可数子覆盖，再构造从属于它的可数指标连续单位分解，便于书写及将局部对象加权求和以作整体构造。但保证这种求和局部只有有限项的是局部有限性，而非指标集的可数性。
\end{remark}

\begin{theorem}
    仿紧 $\land$ $T_2\im$ 任何开覆盖存在从属于它的局部有限的连续单位分解。
\end{theorem}
\begin{proof}
    设开覆盖 $\topo P=\{O_\alpha\}$。由前文，$X$ 为 $T_4$。先用收缩定理，再取局部有限开加细，得到开覆盖 $\{U_i\}$ 及 $\bar U_i\subset O_{\alpha(i)}$。沿用 Smirnov 定理证明中的构造，取闭覆盖 $\{F_i\}$ 使 $F_i\subset U_i$。
    由 Urysohn 引理，取连续函数 $\varphi_i:X\to[0,1]$，在 $F_i$ 上为 $1$，在 $X\backslash U_i$ 上为 $0$，故 $\supp\varphi_i\subset\bar U_i\subset O_{\alpha(i)}$。令
    \[
        \Phi=\sum_i\varphi_i,\qquad
        \rho_\alpha=\frac{\sum_{\alpha(i)=\alpha}\varphi_i}{\Phi}.
    \]
    各和在每点附近仅有有限项，且 $\{F_i\}$ 覆盖 $X$ 保证 $\Phi\geqslant1$。故 $\rho_\alpha$ 连续，$0\leqslant\rho_\alpha\leqslant1$，$\sum_\alpha\rho_\alpha=1$。
    又 $\{\bar U_i\}$ 局部有限，其任意子族之并闭，故
    \[
        \supp\rho_\alpha\subset\bigcup_{\alpha(i)=\alpha}\bar U_i\subset O_\alpha.
    \]
    每点附近只遇到有限个 $\bar U_i$，因而也只遇到有限个 $\supp\rho_\alpha$。即得所求单位分解。
\end{proof}


\section{连通性}

\begin{definition}
    拓扑空间 $X$ 称为\textbf{连通的}(connected)，若它不能写成两个互不相交的非空开集之并。
    定义等价关系如下：$x \sim y$ 当且仅当 $x,y$ 属于 $X$ 的同一个连通子空间。每个等价类称 $X$ 的一个\textbf{连通分支}。
\end{definition}
\begin{theorem}
    连通的两条有用的等价说法：任何一个从 $X$ 到赋予离散拓扑的两点集合 $\{0,1\}$ 的连续映射都不是满射（非平凡开闭集 $A$ 将 $X$ 分成 $A$ 与其补集，而这恰好对应连续满射 $\mathbf1_A:X\to\{0,1\}$）；连通拓扑空间 $X$ 的既开又闭子集只有 $\mt,X$。
\end{theorem}

\begin{theorem}
    连续映射保持连通。
\end{theorem}
\begin{proof}
    置 $A\subset X$ 连通、$f:X\to Y$ 连续，则 $g=f|_A:A\to f[A]$ 为连续满射。
    若 $f[A]=U\cup V$ 为两个不交非空相对开集之并，则 $A=g^{-1}[U]\cup g^{-1}[V]$。由 $g$ 连续且满射，右端为两个不交非空的 $A$ 中开集，与 $A$ 连通矛盾。故 $f[A]$ 连通。
\end{proof}

\begin{definition}
    设区间 $I\subset\R$ 和拓扑空间 $X$，连续映射 $C:I\to X$ 称为 $X$ 上的一条（有向）\textbf{曲线}(curve)。
    $\sigma\in I$ 称为曲线的\textbf{参数}，$\Im C$ 称为\textbf{路径}(path)，不混淆时亦可称曲线。
\end{definition}

\begin{definition}
    设曲线 $C:I\to X$ 及区间 $I'\subset\R$。若存在连续的严格单调双射 $\phi:I'\to I$，则 $C'=C\circ\phi:I'\to X$ 称为 $C$ 的\textbf{重参数化}(reparametrization)。$\phi$ 严格递增时保持方向，严格递减时反转方向；两者均不改变路径，即 $\Im C'=\Im C$。
\end{definition}

\begin{definition}
    拓扑空间 $X$ 称\textbf{路径连通}(path-connected)或\textbf{弧连通}(arcwise-connected)，若 $X$ 中任两点都可被一条完全在 $X$ 中的路径连接。\textbf{局部路径连通}指任点 $x\in X$ 的任意邻域 $N$ 都包含该点的路径连通开邻域 $U\subset N$，即每点都有由路径连通开集组成的邻域基。
\end{definition}
\begin{remark}
    路径连通与局部路径连通无关。要区分“存在一个”与“存在任意小的”：仅仅说存在 $U\ni x$ 还不够，必须要借用 $N$ 来描述任意小。
\end{remark}

\begin{example}
    开球路径连通，因为对任意两点 $a,b\in B(x,r)$，存在线段
 $(1-t)a+tb,t\in[0,1]$ 都在 $B(x,r)$ 里。
 由于这里路径为直线，故路径连通甚至升级为了凸性(convex)。
\end{example}

\begin{theorem}
    连续映射保持路径连通。
\end{theorem}

\begin{proof}
    置 $A\subset X$ 路径连通、$f:X\to Y$ 连续。任取 $y_0,y_1\in f[A]$，则 $\exists a_0,a_1\in A$ 使 $y_0=f(a_0),y_1=f(a_1)$。
    进而存在路径 $\gamma:[0,1]\to A$ 连接 $a_0,a_1$，即 $\gamma(0)=a_0,\gamma(1)=a_1$。复合保持连续性，从而 $f\circ\gamma$ 连接 $y_0,y_1$。故 $f[A]$ 路径连通。
\end{proof}

\begin{example}
    欧氏空间的开集 $O$ 局部路径连通：任给 $x\in O$ 及其在 $O$ 中的邻域 $N$，存在开球 $B(x,r)\subset N$。此球是 $O$ 中路径连通的开邻域。
\end{example}

\begin{theorem}
    设 $A\subset X$ 局部路径连通、$f:X\to Y$ 连续。若限制映射 $g=f|_A:A\to f[A]$ 为开映射，则 $f[A]$ 局部路径连通。
\end{theorem}
\begin{proof}
    任给 $y\in f[A]$ 及其相对开邻域 $O$，取 $x\in A$ 使 $g(x)=y$。因 $g$ 连续，$g^{-1}[O]$ 是 $x$ 在 $A$ 中的开邻域，故存在 $A$ 中路径连通的开集 $U$ 使 $x\in U\subset g^{-1}[O]$。
    由 $g$ 连续且开，$g[U]$ 是 $f[A]$ 中路径连通的开集，且 $y\in g[U]\subset O$。故 $f[A]$ 局部路径连通。
\end{proof}
\begin{remark}
    特别地，局部路径连通空间的连续开像仍局部路径连通。对一般子空间 $A$，不能由 $f$ 开推出 $f|_A$ 开，因为通常只有 $f[A\cap G]\subset f[A]\cap f[G]$。
\end{remark}

\begin{theorem}
    路径连通 $\im$ 连通。
\end{theorem}
\begin{proof}
    置 $X$ 路径连通。反设 $X=U\cup V$，其中 $U,V$ 为不交非空开集。取 $x\in U,y\in V$ 及连接两点的路径 $\gamma:[0,1]\to X$。则 $\gamma^{-1}[U],\gamma^{-1}[V]$ 是 $[0,1]$ 中两个不交非空相对开集，且并为 $[0,1]$，与区间连通矛盾。
\end{proof}

\begin{theorem}
    连通 $\land$ 局部路径连通 $\im$ 路径连通
\end{theorem}
\begin{proof}
    空集情形显然。固定 $p\in X$，记 $C$ 为能以路径与 $p$ 相连的点之集。每点都有路径连通的开邻域 $U$；由路径拼接，$U$ 若与 $C$ 相交便含于 $C$，否则含于 $X\backslash C$。故 $C$ 及其补集均开。因 $p\in C$，连通性给出 $C=X$，故 $X$ 路径连通。
\end{proof}

连通分支未必只有有限个，例如离散空间 $\N$ 的每个单点集都是一个连通分支。

\begin{definition}
    连续映射 $\gamma:[0,1]\to X$ 若满足 $\gamma(0)=\gamma(1)=p$，则称为以 $p$ 为基点的\textbf{闭路}。非空路径连通空间 $X$ 称为\textbf{单连通的}(simply connected)，如果每条闭路都可在 $X$ 内保持基点不动而连续收缩至一点。即存在连续映射 $H:[0,1]^2\to X$，使
    \[
        H(s,0)=\gamma(s),\qquad H(s,1)=p,\qquad H(0,t)=H(1,t)=p.
    \]
\end{definition}

\begin{example}
    $\R^n$ 中的非空凸集单连通：任意两点由线段连接，而闭路 $\gamma$ 可由 $H(s,t)=(1-t)\gamma(s)+tp$ 收缩至基点 $p$。特别地，开球、闭球均单连通。但单连通空间未必同胚于开球，例如正维闭球紧致，开球却不紧致。
\end{example}

同胚于圆周 $S^1$ 的子空间称为\textbf{简单闭曲线}；一般闭路则允许自交。平面上的简单闭曲线将平面分为有界的内部和无界的外部，此为 Jordan 曲线定理。更一般地，Jordan--Brouwer 分隔定理断言：若 $n\geq2$ 且 $\Sigma\subset\R^n$ 同胚于球面 $S^{n-1}=\{x\in\R^n:|x|=1\}$，则 $\R^n\backslash\Sigma$ 恰有两个连通分支，一有界、一无界，且均以 $\Sigma$ 为边界。此处不证。

一般连通区域的边界未必是曲面。物理学中的通量积分与 Gauss 定理还需定向及相应的正则性条件，详见附录 \ref{appx:form}。
